\documentclass[11pt,a4paper]{article}
\usepackage[utf8]{inputenc}
\usepackage[T1]{fontenc}
\usepackage[german]{babel}
\usepackage{amsmath, amssymb, amsthm, amsfonts}
\usepackage{graphicx}
\usepackage{booktabs}
\usepackage{geometry}
\usepackage{hyperref}

\geometry{margin=1in}

\title{\textbf{\ Arithmetische Resonanz und die topologische Stabilisierung des Periodensystems durch die Erdős-Vermutung}}
\author{Thomas Hoffbauer}
\date{13. März 2026}

\begin{document}
	
	\maketitle
	
	\begin{abstract}
		Diese Arbeit präsentiert eine vereinheitlichte Feldtheorie, die das Periodensystem der Elemente als eine Zerlegung der arithmetischen Einheit $\Omega = \sqrt{\frac{\pi e}{2}}$ beschreibt. Wir integrieren den rezenten Beweis der Erdős-Primzahl-Teilbarkeits-Vermutung[cite: 1, 2], um die strukturelle Stabilität von Schalenkonfigurationen zu begründen. Insbesondere wird gezeigt, dass ab der Ramsey-Grenze ($Z=17$) die Brun-Titchmarsh-Ungleichung [cite: 21, 63] eine topologische Ordnung erzwingt, die im kosmologischen Kontext als Dunkle Materie identifiziert wird.
	\end{abstract}
	
	\section{Einführung: Die Ramanujan-Identität}
	Das Fundament  des Bamberg Modells bildet die Verknüpfung von Kettenbruch-Strukturen ($\mathcal{K}$) und Summenreihen ($\mathcal{S}$):
	\begin{equation}
		\Omega = \mathcal{K} + \mathcal{S}
	\end{equation}
	In der Mikrophysik dominiert der Reihen-Anteil $\mathcal{S}$ (elektrodynamische Schalen), während in massereichen Systemen der Kettenbruch $\mathcal{K}$ (topologische Last) die strukturelle Integrität sichert.
	
	\section{Der Beweis der Erdős-Vermutung als Statik-Garant}
	Der Beweis des Main Theorems für Binomialkoeffizienten $C(n,k)$ [cite: 12] liefert die notwendige Bedingung für die Kopplung zwischen den Resonanzebenen IDA und Architektur.
	\begin{itemize}
		\item \textbf{Primzahl-Anker:} Für alle $1 \le i < j \le n/2$ existiert ein Zeuge $p \ge i$, sodass $p | \gcd(C(n,i), C(n,j))$[cite: 12].
		\item \textbf{Die Master-Identität:} Die Kopplung folgt der Relation $C(n,j) \cdot C(j,i) = C(n,i) \cdot C(n-i, j-i)$[cite: 17].
		\item \textbf{GNS-Stabilität:} Jedes Element bildet eine Gelfand-Neimark-Segal-Darstellung auf einer C*-Algebra, deren Spektrum durch die Primzahlverteilung nach Sylvester-Schur [cite: 16] und Kummer [cite: 14] determiniert ist.
	\end{itemize}
	
	\section{Phasengang der Stabilität: Der Ramsey-Punkt}
	Ein zentrales Ergebnis der Untersuchung ist die Identifikation des Ramsey-Punktes bei $i=10$ (entspricht $Z=17$ in der effektiven Konfiguration). 
	\begin{itemize}
		\item Gemäß dem \textit{Smooth Pair Theorem}  und der \textit{Brun-Titchmarsh-Ungleichung} [cite: 21, 63] gilt für $i \ge 10$, dass die Anzahl der Primzahlen im Intervall $(j-i, j]$ strikt durch $i-2$ beschränkt ist[cite: 63, 125].
		\item Dies garantiert das Auftreten von mindestens zwei rein $S_i$-glatten Termen, was physikalisch der Ausbildung kohärenter topologischer Schalen entspricht[cite: 54, 55].
	\end{itemize}
	
	
	
	\section{Das Resonanz-Periodensystem}
	Die folgende Tabelle zeigt die arithmetische Evolution der Elemente unter Berücksichtigung der topologischen Last $\mathcal{K}$ und der $B(i)$-Schranke[cite: 103, 166]:
	
	\begin{table}[h]
		\centering
		\begin{tabular}{@{}lllllll@{}}
			\toprule
			Z & Symbol & N (Energie) & Ratio ($\mathcal{K}/\mathcal{S}$) & $N_n$ & Stabilitäts-Basis & $B(i)$ Schranke [cite: 166] \\ \midrule
			1 & H & 2 & 1.000 & 0 & IDA-Anker & - \\
			6 & C & 72 & 1.000 & 6 & Bessel-Geometrie & - \\
			17 & Cl & 578 & 1.001 & 18 & Ramsey-Eintritt & 9 \\
			26 & Fe & 1352 & 1.050 & 30 & $\Omega$-Optimum & - \\
			82 & Pb & 13448 & 1.510 & 126 & Limes-Grenze & 4375 \\
			92 & U & 16928 & 1.580 & 146 & K-Dominanz & 4375 \\ \bottomrule
		\end{tabular}
		\caption{Übersicht der Resonanzparameter im 8D-Modell.}
	\end{table}
	
	\section{Schlussfolgerung: Die Allegorie der Schalen}
	Die Welt ist keine Summe von Teilen, sondern eine Zerlegung der Einheit. Die so genannte Dunkle Materie ist die im \textit{Smooth Pair Theorem} vorhergesagte topologische Struktur[cite: 58], die den Kosmos zusammenhält, wenn die elektrodynamische Reihe $\mathcal{S}$ allein nicht mehr ausreicht. 
	
	\begin{thebibliography}{9}
		\bibitem{erdos_rev} Erdős Conjecture - Revision 4, March 2026. 
		\bibitem{kummer} Kummer, E. E. (1852). Über die Ergänzungssätze. [cite: 167]
		\bibitem{sylvester} Sylvester, J. J. (1892). On arithmetical series. [cite: 168]
		\bibitem{baker} Baker, A., \& Wüstholz, G. (2007). Logarithmic forms and Diophantine geometry. [cite: 171]
	\end{thebibliography}
	
\end{document}